% Source-derived chapter 05: Stably birational geometry of K-trivial varieties
% Original source: motivic-invariants-birational-maps
\chapter{Stably birational geometry of K-trivial varieties}
\label{motivic-invariants-birational-maps-chapter-05}
\source{motivic-invariants-birational-maps}

\begin{remark}[Source section]
\label{motivic-invariants-birational-maps-chapter-05-source}
\source{motivic-invariants-birational-maps}
\source{motivic-invariants-birational-maps}
This chapter reproduces the corresponding section of the retrieved
LaTeX source, including its definitions, statements, examples, and
proofs. It has not yet been translated into Lean.
\notready
\end{remark}

% The source section title was: Stably birational geometry of K-trivial varieties
\label{app:MRC}
\source{motivic-invariants-birational-maps}
\medskip



The following result is well-known in characteristic zero.
We present the proof 
which works over arbitrary fields as well as a more general statement (see Remark \ref{rem:burt}), both 
communicated to us by Burt Totaro.

\begin{theorem}[Totaro]
\source{motivic-invariants-birational-maps}
\notready\label{thm:BurtCY}
\source{motivic-invariants-birational-maps}
\label{pic1}
\source{motivic-invariants-birational-maps}
Let $X_1$ and $X_2$ be 
smooth projective varieties
over a field with a
birational map
$\phi:X_1 \dto X_2$.
Assume that $K_{X_1}$ and $K_{X_2}$ are nef.
Then $\phi$
is a pseudo-isomorphism.
If in addition 
 $\NS(X_1)_\Q$ 
 or $\NS(X_2)_\Q$
 is one-dimensional, then
$\phi$ is an isomorphism.
\end{theorem}

\begin{proof}
Let us show that $\phi$ is a pseudo-isomorphism, that is both
$\phi$
and $\phi^{-1}$ do not have any exceptional divisors.
We follow the proof of \cite[Corollary 3.54]{KollarMori}, which works
in any characteristic. Let $Y$ be the normalization of the closure
of the graph of $\phi$,
with birational morphisms $g_1\colon Y\to X_1$
and $g_2\colon Y\to X_2$. 
We have a linear equivalence of $\Q$-Weil divisors:
$$K_Y \sim_{\Q} g_i^*(K_{X_i})+Z_i$$
for $i=1,2$, where $Z_i$ is an effective $\Q$-divisor
whose support is the union
of all exceptional divisors of $g_i$ (because $X_i$ are smooth, hence terminal \cite[Claim 2.10.4]{KollarSingMMP}). 
So we have
$$g_1^*(K_{X_1})-g_2^*(K_{X_2})\sim_{\Q} Z_2-Z_1.$$
Here $g_1^*(K_{X_1})$ and $g_2^*(K_{X_2})$ are nef.
Applying the negativity lemma to $g_1 : Y \to X_1$
(resp. $g_2 : Y\to X_2$), we obtain that $Z_2-Z_1$
(resp. $Z_1-Z_2$) is effective \cite[Lemma 3.39]{KollarMori}.
(The negativity lemma holds in any characteristic,
because the proof uses resolution
of singularities only in dimension 2.)
Thus $Z_1=Z_2$ and so $g_1$
and $g_2$ have the same exceptional divisors, hence $\phi$ is a pseudo-isomorphism.


Thus $\phi^*$
gives an isomorphism
\[
\Pic(X_2) \simeq \Pic(X_1),
\]
which takes divisors
to effective divisors (but in general does not preserve ampleness).
Furthermore, 
we have
an induced
isomorphism $\NS(X_2) \simeq \NS(X_1)$ \cite[Example 19.1.6]{Fulton}.
Now assume that
$\NS(X_1)_\Q \simeq \NS(X_2)_\Q \simeq \Q$.
Since ampleness is a numerical condition, in this case
every nonzero
effective divisor is ample, in particular
$\phi^*$ 
takes ample divisors to 
ample divisors.
Take any ample
divisor 
$H_2 \in \Pic(X_2)$ 
and let $H_1 = \phi^*(H_2)$.
For 
every $m \ge 0$,
$\phi$ induces 
an isomorphism $H^0(X_1,\OO(mH_1)) \cong H^0(X_2,\OO(mH_2))$.
These are compatible with products,
that is we have a ring isomorphism
\[
\bigoplus_{m\geq 0} H^0(X_1,\OO(mH_1)) \simeq
\bigoplus_{m\geq 0} H^0(X_2,\OO(mH_2)).
\]
By taking Proj of both sides, it follows that $\phi$ is in fact
an isomorphism from $X_1$ to $X_2$.
\end{proof}

\begin{remark}
\source{motivic-invariants-birational-maps}
\notready\label{rem:burt}
\source{motivic-invariants-birational-maps}
The proof of Theorem~\ref{thm:BurtCY} 
goes through in the relative case: given proper morphisms of varieties over a field $f_1: X_1 \to S$ and $f_2: X_2 \to S$ 
with $X_1$ and $X_2$ smooth (or just $\Q$-Gorenstein
terminal),
assume 
that
$K_{X_i}$ is nef over $S$ for $i = 1, 2$.
Let $\phi: X_1 \dto X_2$
be a birational map over $S$.
Then $\phi$ is a pseudo-isomorphism.
\end{remark}



We need to recall some results about separable maps in positive characteristic.
A dominant map $\phi: X \dto Y$ between $\k$-varieties 
is called separable if the
corresponding field extension $\k(Y)/\k(X)$
is separable~\cite[\href{https://stacks.math.columbia.edu/tag/030I}{Tag 030I}]{stacks-project},
that is $\k(Y)$ is a finite
separable extension of a purely transcendental extension of
$\k(X)$. 
Equivalently,
$\k(X)$ is geometrically
reduced over $\k(Y)$~\cite[\href{https://stacks.math.columbia.edu/tag/05DT}{Tag 05DT}]{stacks-project}.
Recall that a variety $X$ is called separably uniruled
if there exists a 
separable
dominant map
$$Y \times \P^1 \dto X$$
for some variety $Y$ with $\dim Y = \dim X - 1$
(see e.g.~\cite[Definition IV.1.1]{Kollarrat}).
The following lemma should be well-known.

\begin{lem}
\source{motivic-invariants-birational-maps}
\notready\label{lem-defSU}
\source{motivic-invariants-birational-maps}
Let $X$ be a smooth complete variety over an algebraically closed field $\k$.
Suppose that there exist a variety $Y$ and
a separable dominant map
$$u: Y \times \P^1 \dto X$$
such that the proper transform of $\{y\} \times \P^1$ 
is not a point for general $y \in Y$.
Then $X$ is separably uniruled.
\end{lem}

\begin{proof}
By~\cite[Theorem IV.1.9]{Kollarrat},
it suffices to prove the existence of a
morphism
$f : \P^1 \to X$
which is free (this means that $f^*T_X$
is globally generated).
This is shown in~\cite[Lemma 1.2]{MR977778};
for the convenience of the reader, we reproduce the proof.
Since $\k$ is assumed
to be algebraically closed
(hence perfect), after shrinking $Y$, we can assume that $Y$ is smooth.
Therefore $u$ is a morphism outside of a locus of codimension two.
As such a locus does not dominate $Y$,
up to further shrinking $Y$, 
we can assume that $u$ is a morphism.

As $u$ is separable, by generic smoothness~\cite[\href{https://stacks.math.columbia.edu/tag/056V}{Tag 056V}]{stacks-project}
the tangent map $du : T_{Y \times \P^1} \to u^*T_X$
is surjective at a general point $(y,t) \in Y \times \P^1$.
Write $\P^1_y \cnec \{y\} \times \P^1$, then
$(du)|_{\P_y^1}$ is generically surjective.
Thus
we have 
a morphism of sheaves on $\P^1$
\[
\cO_{\P_y^1}(2) \oplus \cO_{\P_y^1}^{\dim Y} \simeq
T_{Y \times \P^1}|_{\P^1_y}  \longrightarrow 
(u^*T_X)|_{\P^1_y} \simeq
\bigoplus_{i = 1}^{\dim X} \cO(a_i)
\]
with torsion cokernel
and it follows immediately
that
all 
$a_i \ge 0$, so 
$(u^*T_X)|_{\P^1_y}$ 
is globally generated 
and $u|_{\P^1_y} : \P^1_y \to X$ is a free morphism.
\end{proof}

Recall that a variety $W$ is called separably rationally connected~\cite[IV.3.2.3]{Kollarrat}
if there exist a variety $S$ and a rational map
$u : S \times \P^1 \dto W$ such that the two-point evaluation map
\begin{equation}\label{eq:rat-conn}
\source{motivic-invariants-birational-maps}
    u_2 = u \times_S u : S \times \P^1 \times \P^1  \dto W \times W
\end{equation}
is dominant and separable.
Note that this condition implies that $u$ is dominant 
and separable.

\begin{lem}
\source{motivic-invariants-birational-maps}
\notready\label{lem-stbirCY} 
\source{motivic-invariants-birational-maps}
Let $Z$ and $Z'$ be smooth projective varieties 
of the same dimension
over 
an algebraically closed field $\k$.
Let $W$ be a separably rationally connected variety over $\k$.
Suppose that both $Z$ and $Z'$ are not separably uniruled.
Then for
every birational map
$\phi : Z \times W \dto Z' \times W$,
there is a unique birational map
$\ol{\phi} : Z \dto Z'$ making the diagram
\begin{equation}\label{eq-ZZW}
\source{motivic-invariants-birational-maps}
    \xymatrix{
Z \times W \ar[d] \ar@{-->}[r]^{\phi} & Z' \times W \ar[d] \\
Z  \ar@{-->}[r]^{\ol{\phi}} & Z'  \\
}
\end{equation}
commutative \textup(here, the vertical arrows are the
projections onto the first factors\textup).
\end{lem}

See~\cite[Theorem 2]{LiuSebag} for a similar result,
where 
a more general statement 
was proven in characteristic zero by Liu and Sebag.
Our proof follows a similar strategy
to that of \cite[Theorem 2]{LiuSebag},
relying on the statement analogous
to \cite[Lemma 5]{LiuSebag}.



\begin{proof}

Let 
$$p: Z \times W \dto Z' \times W \to  Z' $$
be the composition of $\phi$ with the projection.
Since $W$ is separably rationally connected,
there exist a variety $S$ and a rational map
$u : S \times \P^1 \dto W$ such that
the map $u_2$ \eqref{eq:rat-conn}
is dominant and separable.
Since $u$ is also dominant and separable,
so is the composition
$$\Phi: Z \times S \times \P^1 
\overset{\Id_Z \times u}{\dto} Z  \times W
\overset{p}{\dto} Z'.$$

Since $Z'$ is not separably uniruled, $\Phi$ contracts
$\{(z, s)\} \times \P^1$ for general $(z, s) \in Z \times S$.
In other words, $p(z,-): W \dto Z'$ 
contracts general rational curves of $W$ parameterized
by $S$ through $u$. 
Since these curves connect general points of $W$
because $u_2$ is dominant,
$p(z,-)$ contracts $W$ for general $z \in Z$. 

Let $\Gamma \subset (Z \times W) \times Z'$ be the closure of the graph
of $p$ and let $\ol{\Gamma} \subset Z \times Z'$ 
be the image of $\Gamma$.
The fiber of $\ol{\Gamma} \to Z$ over a general point 
$z \in Z$ is the point
$p(\{z\} \times W)$. 
Moreover, $\ol{\Gamma} \to Z$ has degree one,
because the composition of projections
$$(Z \times W) \times Z' \to Z \times Z' \to Z$$
induces the composition of surjective morphisms
$$\Gamma_w \cnec \Gamma \cap (Z \times \{w\} \times Z')
\to \ol{\Gamma} \to Z,$$
which is birational for general $w \in W$;
this is because $\Gamma_w$ is the graph of the rational
map $p|_{Z \times \{w\}}$.
Therefore $\ol{\Gamma}$
is the graph of a rational dominant map $\ol{\phi}: Z \dto Z'$,
and it makes~\eqref{eq-ZZW} commutative by construction.
It is clear from~\eqref{eq-ZZW} that $\phi$ uniquely determines $\ol{\phi}$.

To show that $\ol{\phi}$ is birational,
we apply the same construction
to $\phi^{-1}$, and by uniqueness
both compositions of
$\ol{\phi}$ with
$\ol{\phi^{-1}}$ are identity maps.
\end{proof}




\begin{cor}
\source{motivic-invariants-birational-maps}
\notready\label{cor-CYnSB} 
\source{motivic-invariants-birational-maps}
Let $Z$ and $Z'$ be K-trivial varieties 
of Picard number $1$
over 
an algebraically closed field $\k$.
If $Z \not\simeq Z'$,
then $Z \times W$ is not birational to $Z' \times W$
for any separably rationally connected variety $W$ over $\k$.
In particular, $Z$ is not stably birational to $Z'$.
\end{cor}




\begin{proof} 

Since K-trivial varieties are not separably uniruled~\cite[Corollary IV.1.11]{Kollarrat}, 
Theorem~\ref{thm:BurtCY}
and Lemma~\ref{lem-stbirCY} show that
$Z \times W$ is not birational to $Z' \times W$.
\end{proof}
